\section{Periodic points of a special K\"ahler group}\label{sec-period}
In this section, we study the periodic points of an automorphism group, as a linear quotient of some special K\"ahler group, with 
the main results {\hyperref[prop-sol-fixed]{Proposition    \ref*{prop-sol-fixed}}} and {\hyperref[prop_fund_pp]{Proposition     \ref*{prop_fund_pp}}}.
As an application, we apply {\hyperref[prop_fund_pp]{Proposition     \ref*{prop_fund_pp}}}  to the locally constant MRC fibration $f:X\to Y$ in \cite[Theorem 1.2]{MW21}, which plays a significant role to find a section of such $f$  (cf.~{\hyperref[cor-lc-periodic]{Corollary \ref*{cor-lc-periodic}}} and  {\hyperref[lemma_fixed-pt-subbundle]{Lemma   \ref*{lemma_fixed-pt-subbundle}}}). 
We refer readers to \cite[Section 4]{LOY19} which considered the Albanese map onto the complex torus; however, the fundamental group of our $Y$ (as the object of the MRC fibration) here is neither abelian nor solvable.



To begin with this section,  we briefly recall some basic facts about the Shafarevich map, which will be used in the proof of this and also later sections. 
\begin{defn}[{\cite[Definition 3.5]{Kollar95}}]\label{defn-shafarevich}
Let $Y$ be a normal variety and $H\lhd \pi_1(Y)$ a normal subgroup of the fundamental group of $Y$.
A normal variety $\textup{Sh}^H(Y)$ and a rational map $\textup{sh}_Y^H:Y\dashrightarrow\textup{Sh}^H(Y)$ are called the \textit{$H$-Shafarevich variety} and the \textit{$H$-Shafarevich map} of $Y$ if
\begin{enumerate}
    \item[(1)] $\textup{sh}_Y^H$ has connected fibers, and
    \item[(2)] there are countably many closed subvarieties $D_i\subseteq Y$ $(D_i\neq Y)$ such that for every closed, irreducible subvariety $Z\subseteq Y$ with $Z\not\subseteq\cup D_i$, we have 
    $$\textup{sh}_Y^H(\widetilde{Z})=\textup{point}~~\textup{if and only if}~~\textup{im}[\pi_1(\widetilde{Z})\to\pi_1(X)]\textup{~has finite index in }H,$$
where $\widetilde{Z}$ is the normalization of $Z$.
\end{enumerate}
\end{defn}

The following proposition guarantees the existence of the Sharfarevich maps.
\begin{prop}[{cf.~\cite[Theorem 3.6]{Kollar95}}]\label{prop-kol-sha}
Let $X$ be a normal variety, and $H\lhd\pi_1(X)$ a normal subgroup. Then
\begin{enumerate}
\item[(1)]The $H$-Shafarevich map $\textup{sh}_X^H:X\dashrightarrow\textup{Sh}^H(X)$ exists.
\item[(2)]For every choice of $\textup{Sh}^H(X)$ (within its birational equivalence class) there are open subsets $X_\circ\subseteq X$ and $W_\circ\subseteq\textup{Sh}^H(X)$ with the following properties:
\begin{itemize}
    \item $\textup{sh}_X^H:X_\circ\to W_\circ$ is everywhere defined on $X_\circ$.
    \item Every fibre of $\textup{sh}_X^H|_{X_\circ}$ is closed in $X$.
    \item $\textup{sh}_X^H|_{X_\circ}$ is a topologically locally trivial fibration.
\end{itemize}
\item[(3)] If $X$ is proper, then $\textup{sh}_X^H:X_\circ\to W_\circ$ is proper and the very general points of $X$ is contained in $X_\circ$ for a suitable choice of $\textup{Sh}^H(X)$ and $X_\circ$.
\end{enumerate}
\end{prop}

Now, we apply {\hyperref[defn-shafarevich]{Definition   \ref*{defn-shafarevich}}}  and {\hyperref[prop-kol-sha]{Proposition  \ref*{prop-kol-sha}}} to our specific setting. 
Let $Y$ be a normal projective variety and let $\chi:\pi_1(Y)\rightarrow \text{GL}(r,\mathbb{C})$ be a representation. 
%By Selberg's lemma (\cite{Sel60}), every finitely generated linear group over $\CC$ is virtually torsion-free, and consequently there exists a finite \'etale cover 
Let $\pi:Y'\rightarrow Y$ be a finite \'etale cover such that the algebraic Zariski closure $\overline{G_{Y'}}$ of the image $G_{Y'}:=\chi(\pi_*(\pi_1(Y')))$ of the following composed map
\[
\chi':\pi_1(Y') \xrightarrow{\pi_*} \pi_1(Y) \xrightarrow{\chi} \text{GL}(r,\CC)
\]
is connected. 
Moreover, denote by $\text{Rad}(\overline{G_{Y'}})\lhd \overline{G_{Y'}}$ for the solvable radical of $\overline{G_{Y'}}$. 
Then the quotient $\overline{G_{Y'}}/\text{Rad}(\overline{G_{Y'}})$ is torsion-free. 
Denote by $K\lhd \pi_1(Y')$ the kernel of the map $\pi_1(Y')\rightarrow \overline{G_{Y'}}/\text{Rad}(\overline{G_{Y'}})$, which is normal in $\pi_1(Y')$.  
Then it follows from {\hyperref[prop-kol-sha]{Proposition  \ref*{prop-kol-sha}}} that there exists a dominant almost holomorphic map (which is the $K$-Shafarevich map) $\text{sh}_{Y'}^K:Y'\dashrightarrow \text{Sh}^K(Y')$ to a smooth projective variety. 
In other words, there exists a dense Zariski open subset $Y'_{\circ}\subset Y'$ such that the restriction $\text{sh}^K_{Y'}|_{Y'_{\circ}}$ is well-defined and proper. %This map $\text{sh}_{Y'}^K$ is called the \emph{$K$-Shafarevich map} of $Y'$ (cf.~\cite[Definition 3.5]{Kollar95}). In particular, we borrow the following result which sits in a more general setting:
Take  a very general point $x\in Y'_{\circ}$ and let $Z$ be a subvariety through $x$, with normalisation $n:\widetilde{Z}\rightarrow Z$.   
Then, the rational map $\text{sh}^K_{Y'}$ maps $Z$ to a point if and only if the composed map
\[
\pi_1(\widetilde{Z})\xrightarrow{n_*} \pi_1(Y') \rightarrow \overline{G_{Y'}}\rightarrow \overline{G_{Y'}}/\text{Rad}(\overline{G_{Y'}})
\]
has finite image (cf.~{\hyperref[defn-shafarevich]{Definition   \ref*{defn-shafarevich}}}).  
We refer readers to \cite{CCE15} for more information on the Shafarviech maps and Shafarviech varieties. 

The following theorem is crucial for our proof of {\hyperref[t.mainthm]{Theorem \ref*{t.mainthm}}}.

\begin{thm}[{\cite[Theorem 1]{CCE15}}]
	\label{t.basemfdShafarevichmap}
With the notation above, assume that $Y'$ has at worst klt singularities. Then the smooth projective variety $\text{Sh}^K(Y')$ is of general type.
\end{thm}

\begin{proof}
	Let $\eta:\widetilde{Y}'\rightarrow Y'$ be a resolution. By \cite[Theorem 1.1]{Tak03}, the push-forward $\eta_*:\pi_1(\widetilde{Y}')\rightarrow \pi_1(Y')$ is an isomorphism. In particular, the composed map 
	\[
	\widetilde{Y}' \rightarrow Y' \dashrightarrow \text{Sh}^K(Y')
	\]
	is a Shafarevich map for $K\lhd \pi_1(\widetilde{Y}')$. Then the statement follows immediately from \cite[Th\'eor\`eme 1]{CCE15} and the discussion above.
\end{proof}


In the sequel of this section, we study the periodic point of certain groups acting on a projective varieties. By the Borel fixed point theorem, we generalize \cite[Theorem 4.1]{LOY19} to the solvable group case to get the first main result of this section.
Compared with \cite[Theorem 4.1]{LOY19}, our argument is in the line of the theory of algebraic group.


Let $\sigma:G\times F\to F$ be a group action of $G$ on a projective variety $F$.
We say that $G$ has a \textit{fixed point} $y\in F$, if  $\sigma(g,y)=y$ for any $g\in G$.
We say that $G$ has a \textit{periodic point} $y\in F$, if there exists a positive integer $m$ such that $\sigma(g^m,y)=y$ for any $g\in G$. 

\begin{prop}\label{prop-sol-fixed}
	Let $G$ be a group acting on a projective variety $F$ via a group homomorphism $\rho: G\to \textup{Aut}(F)$ to the automorphism group of $F$ such that the image $\rho(G)$ is virtually solvable and there is a $G$-linearized ample line bundle $L$ on $F$.
	Then $G$ has a periodic point on $F$.	
\end{prop}
\begin{proof}
	After replacing  $G$ by a finite index subgroup, we can assume that $\Gamma:=\rho(G)$ is solvable. By hypothesis, $\rho(G)$ is a subgroup of the linear algebraic group $\textup{Aut}_L(F)$.
	Denote by $[\Gamma,\Gamma]$ the commutator subgroup of $\Gamma$.  
	Recall that for an algebraic group $H$, its commutator group $[H,H]$ is a closed subgroup of $H$, hence also algebraic by \cite[Chapter I, \S 2.3 Proposition, pp.~58-59]{Bor69}.
	Hence,  %$\overline{[\Gamma,\Gamma]}\subseteq [\overline{\Gamma},\overline{\Gamma}]$ and thus
	we have $\overline{[\Gamma,\Gamma]}=[\overline{\Gamma},\overline{\Gamma}]$ by \cite[Chapter I, \S2.1(e), p.~57]{Bor69}.	
	Inductively, the closure of the $p$-th derived group  $\overline{\Gamma^{(p)}}$ coincides with the $p$-th derived group $\overline{\Gamma}^{(p)}$ of $\overline{\Gamma}$.
	
	By the solvability of $\Gamma$,  the $p$-th derived group $\Gamma^{(p)}$ is trivial for some $p<\infty$ (i.e., the derived length of $\Gamma$ is finite).
	Hence, $\overline{\Gamma}^{(p)}$ is also trivial for some $p<\infty$ and thus $\overline{\Gamma}$ is a solvable algebraic group.
	With 
	%$\overline{\Gamma}$ 
	$G$ replaced by some subgroup of finite index if necessary, we can  assume that $\overline{\Gamma}$ is connected.
	By Borel's fixed-point theorem (cf.~\cite[Theorem 10.4, p.~137]{Bor69}), the action of $\overline{\Gamma}$ on $F$ has a fixed point and our proposition is proved by considering the exact sequence $1\to\ker\to G\to G|_F\to 1$.
\end{proof}


In what follows, we recall the notion of {\it special varieties} in the sense of Campana. 
Roughly speaking, a special variety is a compact complex variety in the Fujiki class $\mathcal{C}$ (i.e., bimeromorphic to a compact K\"ahler manifold) that admits no orbifold-theoretic (in the sense of F. Campana) meromorphic fibration onto a positive-dimensional variety of (orbifold) general type. For the precise definition, see \cite[Definition 2.1(2)]{Cam04}. 

%we refer readers to \cite[\S 2]{Cam04} for the definition of \textit{special} (analytic) varieties in the sense of Campana. 
Note that there are two fundamental examples of special varieties (\cite[Theorem 3.22, Theorem 5.1]{Cam04}) as indicated in the following lemma; in fact, every special variety is essentially built up from them (\cite[Example 2.3, \S 6.5]{Cam04}).

\begin{lemme}\label{lem-kappa-special}
Let $X$ be 
a compact complex variety in the Fujiki class $\mathcal{C}$.
%a compact K\"ahler manifold. 
Then $X$ is special if one of the following holds:
\begin{itemize}
\item[\rm(1)] $X$ is rationally connected;
\item[\rm(2)] The Kodaira dimension vanishes, i.e., $\kappa(X)=0$.	
\end{itemize}
\end{lemme}

\begin{rmq}
One should notice that the notion `rational connectedness' in \cite{Cam04} (cf. \cite[\S 3.3]{Cam04}) is usually called `rational chain connectedness' (cf. \cite[3.2 Definition]{Kollar96}), and this is why in \cite[Theorem 3.22]{Cam04} the smoothness of $X$ is required. In this article, we take the usual definition of `rational connectedness' (i.e., any two general points can be connected by a rational curve), and since this is by definition a bimeromorphic property for complex varieties, we do not need the smoothness condition in the lemma. Let us remark that `rational chain connectedness' is however not a bimeromorphic property, this is why the lemma does not hold in this generality; yet the two notions `rational connectedness' and `rational chain connectedness' coincide for varieties with dlt singularities by \cite[Corollary 1.5(2)]{HM07}.      
\end{rmq}

Besides, special varieties have the following property, known as the `weak speciality' (cf. \cite[\S 9]{Cam04}):
%Besides, the following lemma is known to experts.
\begin{lemme}[{cf.~\cite[Proposition 9.27]{Cam04}}]\label{lem-special-no-to-gt}
Let $X$ be a special variety. 
%lying in the Fujiki class $\mathcal{C}$ (i.e., a special variety bimeromorphic to a compact K\"ahler manifold). 
Then any finite \'etale cover of $X$ can never dominate a positive dimensional variety of general type.
\end{lemme}




%The above property is also said to be \textit{weak speciality} (cf.~\cite[Section 9]{Cam04}).  
In the sequel we will prove the main results of this section concerning the special K\"ahler groups. First recall that a group $G$ is said to be a \textit{K\"ahler group} if $G$ can be realized as the fundamental group of a compact K\"ahler manifold. 
Further, a K\"ahler group $G$ is said to be \textit{special}, if $G$ can be realized as the fundamental group of a special compact K\"ahler manifold.

Applying \cite{Tak03} to the resolution of some projective klt pair, we immediately have the following lemma.
\begin{lemme}\label{lem-special-klt}
Let $(X,\Delta)$ be a projective klt pair.
If $X$ is special, then the fundamental group $\pi_1(X)$ is a special K\"ahler group.
\end{lemme}


\begin{prop}\label{prop-special-Kahler-abe}
Any linear quotient of a special K\"ahler group is virtually abelian.
\end{prop}

%Before we prove {\hyperref[prop-special-Kahler-abe]{Proposition   \ref*{prop-special-Kahler-abe}}}, we refer readers to 



%Now we come back to the proof of {\hyperref[prop-special-Kahler-abe]{Proposition   \ref*{prop-special-Kahler-abe}}}. 
\begin{proof}%[Proof of {\hyperref[prop-special-Kahler-abe]{Proposition   \ref*{prop-special-Kahler-abe}}}]
Let $X$ be a compact K\"ahler manifold and $G:=\pi_1(X)$ the fundamental group.
Let $\chi:G\to \textup{GL}(r,\mathbb{C})$ be a linear representation.
Let $G_X:=\chi(G)$ be the image and $\overline{G_X}$ its Zariski closure. 
We will prove that $G_X$ is virtually abelian. To this end, we can freely replace $G$ (and thus $G_X$ and $\overline{G_X}$) by a finite index subgroup, and correspondingly replace $X$ by a finite \'etale cover, 
%After replacing $X$ by a finite \'etale cover, 
so that we may assume that $\overline{G_X}$ is a connected algebraic group.
Denote by $\overline{R}$ the solvable radical of $\overline{G_X}$ and let $R:=\overline{R}\cap G_X$.
Let $s:G_X\to G_X/R$ be its natural quotient.

By \cite[Th\'eor\`eme 6.5 and its proof]{CCE15}, after replacing $X$ by a further  \'etale cover, there is a proper modification $X'\to X$ such that $X'$ dominates a variety  which is bimeromorphic to the total space of a smooth fibration $\textup{Sh}^{K_1}(X)\to W$ over the Shafarviech variety $\textup{Sh}^{K_2}(X)$ such that  $W$ is of general type, where $K_1$ is the kernel of $\chi$ and $K_2$ is the kernel of $s\circ\chi$. 
Since $X$ is special, our $W$ has to be a single point (cf.~{\hyperref[lem-special-no-to-gt]{Lemma  \ref*{lem-special-no-to-gt}}}).
This in turn implies that the Shafarviech variety $\textup{Sh}^{K_2}(X)$  is a single point and thus the image $s(G_X)=s\circ\chi(\pi_1(X))$ is finite.
Hence, $G_X$ itself is virtually solvable (also cf.~\cite[Th\'eor\`eme 6.3]{CCE15}).
Now, it follows from \cite[Corollaire 4.2]{CCE15} that $G_X$ is virtually abelian, which gives our proposition.
\end{proof}



The following corollary is an immediately consequence of {\hyperref[prop-special-Kahler-abe]{Proposition  \ref*{prop-special-Kahler-abe}}}. 
%We refer readers to \cite[Section 11]{GGK19} for the linear representations of the fundamental group of a klt projective variety with the numerical trivial canonical divisor and the vanishing augmented irregularity.

\begin{cor}\label{cor-virtually-abelian}
Let $X$ be a normal projective variety with at worst klt singularities.
Suppose that $X$ is special in the sense of Campana (this is the case when $\kappa(X)=0$; cf.~{\hyperref[lem-kappa-special]{Lemma  \ref*{lem-kappa-special} (2)}}).
Then any linear quotient of $\pi_1(X)$ is virtually abelian.
\end{cor}



\begin{proof}
By {\hyperref[lem-special-klt]{Lemma \ref*{lem-special-klt}}}, 
$\pi_1(X)$ is a special K\"ahler group (cf.~\cite{Tak03}),
noting that the specialness is a birational invariant. 
Then our corollary follows from {\hyperref[prop-special-Kahler-abe]{Proposition  \ref*{prop-special-Kahler-abe}}}. 
\end{proof}

Now we are in the position to prove the second main result in this section.

\begin{prop}\label{prop_fund_pp}
Let $G$ be a special K\"ahler group.
Suppose that  $G$ acts on a projective variety $F$ via a group homomorphism $\rho:G\to\Aut(F)$ to the automorphism group of $F$, such that there exists a $G$-linearized  ample line bundle $L$ on $F$  (cf. \cite[Defintion 3.2.3, Lemma 3.2.4]{Brion18}).
Then $G$ has a periodic point on $F$.
\end{prop}






\begin{proof}
	It follows directly from {\hyperref[prop-sol-fixed]{Proposition \ref*{prop-sol-fixed}}} and {\hyperref[prop-special-Kahler-abe]{Proposition \ref*{prop-special-Kahler-abe}}}.
\end{proof}


Note that the solvability of special K\"ahler groups has its own  interests (cf.~\cite[Conjecture 7.1]{Cam04}). Thus we may ask the following question.

\begin{ques}\label{ques-kernel-sol}
	Let $G$ be a special K\"ahler group.
	Suppose that $G$ acts on a projective variety $F$ via a group homomorphism $G\to\Aut(F)$ such that there is a $G$-linearized  ample line bundle $L$ on $F$.
	Can we give some descriptions on the kernel of $\rho: G\to G|_F$? 
	Will $\ker\rho$ (and hence $G$) be virtually solvable?
\end{ques}

 
In our case, the solvability of $G$ follows from the solvability of $\ker\rho$ due to the exact sequence $1\to \ker\rho\to G\to G|_F\to 1$  (cf.~{\hyperref[prop-special-Kahler-abe]{Proposition  \ref*{prop-special-Kahler-abe}}}). 
Moreover, {\hyperref[ques-kernel-sol]{Question   \ref*{ques-kernel-sol}}} has a negative answer if we don't assume
the speciality.  
For example, 
it is known that the group given by the presentation
$$\Gamma_g=\left<\alpha_1,\cdots,\alpha_g,\beta_1,\cdots,\beta_g:\prod_{i=1}^g[\alpha_i,\beta_i]=1\right>$$
is K\"ahler.
Indeed, it is the fundamental group of a compact Riemann surface of genus $g$.
However, when $g>1$ the commutator subgroup $[\Gamma_g,\Gamma_g]$ is a free non-abelian group and hence $\Gamma_g$ is not virtually solvable. 
%(cf.~https://arxiv.org/pdf/2101.05905.pdf, Theorem B)?

%\begin{rmq}
%Indeed, the existence of the ample $G$-invariant line bundle has its own interest. 
%In this situation, $G|_F$ would be virtually contained in the identity component $\textup{Aut}_0(F)$ of $\textup{Aut}(F)$ by Fujiki and Lieberman (cf.~\cite{Fuj78} and \cite{Lie78}).
%However, little things were know whether $G$ is solvable or not (cf.~Question \ref{ques-kernel-sol}).
%\end{rmq}

Applying {\hyperref[prop_fund_pp]{Proposition  \ref*{prop_fund_pp}}} to a locally constant fibration, we end up this section with the following corollary. % which will be used in the proof of {\hyperref[mainthm-section]{Theorem \ref*{mainthm-section}}}.
\begin{cor}\label{cor-lc-periodic}
Let $f:X\to Y$ be a locally constant fibration with a fibre $F$.
Suppose that the irregularity $q(F)=0$ and $Y$ is special.
Then the fundamental group $G:=\pi_1(Y)$ acts on $F$ and has a periodic point on it.
\end{cor}


\begin{proof}
To apply {\hyperref[prop_fund_pp]{Proposition  \ref*{prop_fund_pp}}}, we only need to verify that there is a $G$-linearized ample line bundle $L$ on $F$ (cf.~\cite[Definition 2.6]{MW21}). 
Since $f$ is locally constant, we have the following diagram
\[
\xymatrix{F\times Y^{\textup{univ}}\ar[r]\ar[d]&Y^{\text{univ}}\ar[d]\\
X\ar[r]&Y}
\]
where $Y^{\textup{univ}}$ is the universal cover of $Y$ and $F\times Y^{\textup{univ}}\cong X\times_Y Y^{\textup{univ}}$ with the induced projection $\textup{pr}_1: F\times Y^{\textup{univ}}\to F$ and  $\textup{pr}_2: F\times Y^{\textup{univ}}\to Y^{\textup{univ}}$.
Let us fix an $f$-ample divisor $A$ on $X$ and denote by $p_X: F\times Y^{\textup{univ}}\to X$ the natural projection.
Applying \cite[Proof of Lemma 2.7]{MW21}, we see that there exists a line bundle $A_Y$ on $Y$ such that 
$$p_X^*(A-f^*A_Y)\cong\textup{pr}_1^*A_F$$ 
for some line bundle $A_F$ on $F$.
Since $X\cong (F\times Y^{\textup{univ}})/G$, our $p_X^*(A-f^*A_Y)$ is $G$-linearized, and hence,  
$$g^*p_X^*(A-f^*A_Y)\sim p_X^*(A-f^*A_Y)$$ 
for any $g\in G$ (regarded as an automorphism of $F$); see \cite[Defintion 3.2.3, Lemma 3.2.4]{Brion18}.  
Hence, $\textup{pr}_1^*A_F$ is also $G$-linearized.
Moreover, since $G$ acts on $F\times Y^{\textup{univ}}$ via the diagonal, our $A_F$ is a $G$-linearized ample line bundle, noting that $A_F\cong A|_F$.
\end{proof}




\vskip 2\baselineskip





















%{\color{blue} The following is a discussion on $G_1$, and can be removed finally.}
%Let $X\to Y$ be a locally constant fibration. 
%We denote by $G_1:=\{g\in G~|~g|_F=\text{id}\}$ and consider the following diagram:
%\[
%\xymatrix{F\times Y^{\textup{univ}}\ar[r]\ar[d]&Y^{\text{univ}}\ar[d]\\
%F\times (Y^{\text{univ}}/G_1)\ar[r]\ar[d]&Y^{\text{univ}}/G_1\ar[d]\\
%X\ar[r]&Y}
%\]



%From \cite[Theorem 6.5]{CCE15}, we know that:
%Let $Y$ be a compact K\"ahler manifold and $\rho:\pi_1(Y)\to \textup{GL}_N(\mathbb{C})$ a linear representation of its fundamental group.
%Then after replacing $Y$ by a finite \'etale cover, the Shafarevich manifold $\textup{Sh}_\rho(Z)$ associated with $\rho$ is (bimeromorphic to) a compact K\"ahler manifold $Z$ which is the total space of a smooth torus fibration over a manifold of general type.
%$$Y\xrightarrow{\textup{sh}_\rho}Z:=\textup{Sh}_\rho(Y)\xrightarrow{s_\rho}S_\rho(Y)$$
%The manifold $Z$ is bimeromorphic to the Shafarevich manifold of $Y$, obtained from $Y$ by contracting the submanifolds, the fundamental groups of which lie in the kernel of $\rho$.
%The torus fibration (i.e., the submersion $Z\to S_\rho(Y)$) has the tori fibre being the maximal manifolds, on the fundamental groups of which, $\rho$ has an abelian image in $\textup{GL}_N(\mathbb{C})$.
%Note that $\rho$ factors as $\pi_1(Y)\to\pi_1(Z)\to \rho(\pi_1(Y))$ (cf.~\cite[Section 2]{CCE15}).

